\begin{landscape}
\definecolor{color_30879}{rgb}{0,0.12549,0.376471}
\pagenumbering{gobble} % Seitenzahlen entfernen
    % top right corner
    \thispagestyle{fancy}
    \fancyhf{}
    \rhead{\includegraphics[width=0.4\textwidth]{resources/Logo.png}}

\section*{Prozessdokumentation}
\textbf{Für die Nutzung von Künstlicher Intelligenz (KI) in wissenschafftlichen Arbeiten}
\vspace{3mm}

\noindent\textbf{\large Eziz Muhammedkulyyev}
\vspace{3mm}

\begin{longtable}{|p{6.5cm}|p{5cm}|p{8cm}|p{4cm}|}
\hline
\rowcolor{color_30879}
\textbf{\textcolor{white}{Anfrage an folgende KI (inkl. Abrufdatum)}} &
\textbf{\textcolor{white}{Grund für die Verwendung der KI}} &
\textbf{\textcolor{white}{Bewertung der Qualität des KI-Outputs}} &
\textbf{\textcolor{white}{Verwendete Software und Versionsnummer}} \\
\hline
\endhead
\hline
„Warum werden beim direkten XES-Import von \texttt{DomesticDeclarations.xes} und \texttt{InternationalDeclarations.xes} in Celonis keine Case-Attribute wie \texttt{Amount} oder \texttt{Permit ID} als eigene Case-Tabelle angelegt, und wie lässt sich das in Celonis prüfen?“ (16.10.2026) &
Klärung technischer Fragen zu Celonis und zu den Einschränkungen des nativen XES-Imports (Kapitel 4.3) &
Die Erklärung, dass der Standard-XES-Import in Celonis nur die Event-Ebene in eine Aktivitätstabelle übernimmt und Trace-Attribute abflacht, war zutreffend. Ich habe dies im Celonis-Datenmodell überprüft und anhand der offiziellen Celonis-Dokumentation sowie des Community-Eintrags (Buijs, 2019) verifiziert. &
Google Gemini / Antigravity \\
\hline
„Erstelle PQL-Formeln in Celonis für die Durchlaufzeit zwischen \texttt{Declaration SUBMITTED by EMPLOYEE} und \texttt{Payment Handled}, für Mehrfacheinreichungen sowie Filter für abgelehnte Cases.“ (mehrfach, 15.--17.10.2026) &
Teilweise Generierung und Fehlersuche bei PQL-Abfragen für eigene Analyse-Views (\textit{TPT Analysis}, \textit{Rejections Dashboard}) und KPIs (Kapitel 4 und 5.3) &
Die vorgeschlagenen PQL-Ausdrücke (u.\,a.\ mit \texttt{CALC\_THROUGHPUT}, \texttt{FIRST\_OCCURRENCE} und \texttt{PU\_COUNT}) waren als Vorlage hilfreich, mussten jedoch teilweise an die exakten Spaltennamen im Celonis-Datenmodell angepasst werden. Alle Abfragen habe ich in Celonis selbst ausgeführt und die Ergebnisse an Einzelfällen geprüft. &
Google Gemini / Antigravity \\
\hline
„Wie unterscheiden sich die ausgewiesenen Durchlaufzeiten in den Celonis-Standard-Views (\textit{Process Overview} vs.\ \textit{Throughput Times}) und wie verhält sich der Conformance-Checker beim Erweitern des Referenzmodells auf die Top-3-Varianten?“ (16.10.2026) &
Rückfragen zur Funktionsweise von Celonis bei der Kennzahlenaggregation und beim Conformance Checking (Kapitel 4.1 und 4.5) &
Die Hinweise zur internen Ausreißer-Filterung im \textit{Process Overview} und zur Generierung von XOR-Gateways im Conformance-Checker waren nachvollziehbar. Die fachliche Interpretation (insbesondere zum Rollout-Effekt 2018 und zur Rolle \texttt{REJECTED by EMPLOYEE}) habe ich selbst anhand der Dashboards und des \textit{Case Explorers} erarbeitet. &
Google Gemini / Antigravity \\
\hline
„Formuliere aus meinen Celonis-Kennzahlen, Screenshots und Stichpunkten zu Durchlaufzeiten, Bottlenecks, Ablehnungen, Conformance und KPIs einen wissenschaftlich strukturierten LaTeX-Text für Kapitel 1, 4 und 5.3.“ (mehrfach, 16.--18.10.2026) &
Textgenerierung und sprachliche Ausformulierung in LaTeX auf Basis eigener Analyseergebnisse und Notizen &
Die generierten Textabschnitte setzten meine vorgegebenen Kennzahlen, Beobachtungen und Gliederungspunkte sprachlich sauber in LaTeX um. Sämtliche Zahlenwerte, Aussagen, Tabellen und Literaturverweise wurden von mir kritisch gegengelesen, teilweise umformuliert und mit meinen Celonis-Auswertungen abgeglichen. &
Google Antigravity / Gemini \\
\hline
\end{longtable}

\newpage

\noindent\textbf{\large Jesco Cedric Homberg}
\vspace{3mm}

\begin{longtable}{|p{6.5cm}|p{5cm}|p{8cm}|p{4cm}|}
\hline
\rowcolor{color_30879}
\textbf{\textcolor{white}{Anfrage an folgende KI}} &
\textbf{\textcolor{white}{Grund für die Verwendung der KI}} &
\textbf{\textcolor{white}{Bewertung der Qualität des KI-Outputs}} &
\textbf{\textcolor{white}{Verwendete Software und Versionsnummer}} \\
\hline
\endhead
\hline
Zeilenweise Code-Vervollständigung während des Schreibens in der IDE (laufend) &
Beschleunigung wiederkehrender Schreibarbeit im Code &
Vorschläge ohne inhaltliche Tragweite. Jede Übernahme erfolgte bewusst und wurde beim Ausführen der Notebooks mitgeprüft. &
GitHub Copilot (Code-Completion) in PyCharm \\
\hline
Erklärung von Funktionalitäten von Pandas und PM4Py (laufend) &
Tiefenwissenslücken ergänzen &
Vorschläge ohne inhaltliche Tragweite. Es ging um Übersicht und Möglichkeiten, um zu prüfen, was nachgebaut werden muss und was bereits besteht. &
GitHub Copilot (Code-Completion) in PyCharm \\
\hline
Umsetzung eigener Entwürfe und Ablaufskizzen in ausführbaren Python-Code (laufend) &
Werkzeug zur Code-Generierung anhand des eigenen Entwurfs, da das eigene Python-Wissen breit, aber nicht in der nötigen Tiefe vorhanden ist &
Aufbau, Reihenfolge und inhaltliche Festlegungen stammen von mir, die KI hat sie in Code übersetzt. Vorschläge wurden ausnahmslos als lokaler Git-Diff gesichtet und erst nach eigener Prüfung übernommen. &
Claude (Anthropic), Opus 4.8 (GitHub Copilot, Max/1M) \\
\hline
Umsetzung eigener Entwürfe und Ablaufskizzen in ausführbaren Python-Code (Bsp. 1) &
Funktionsgenerierung &
Prompt: Bau mir bitte in dieses Notebook einen simplen Parser für die XES-Dateien unter resources (objektorientiertes Ziel). Cool wäre das in einer SQLite-Datenbank direkt zu speichern, damit ich direkt die strukturiert untersuchen kann zbd das nicht mehrfach machen muss. Ich glaube Events haben immer die gleichen 5 Felder, aber Trace scheint wohl teilweise auch mehr zu haben. Vielleicht macht da dann auch Sinn, die Nicht-Standardfelder in einer Map in einer Spalte AdditionalInfos oder sowas zu hinterlegen. Ein Informationsverlust darf nciht geschehen bei der Transformation zur DB &
Claude (Anthropic), Haiku 4.5 03.10.26 // 6.8 Credits (GitHub Copilot, Max/1M) \\
\hline
Umsetzung eigener Entwürfe und Ablaufskizzen in ausführbaren Python-Code (Bsp. 2) &
Funktionsgenerierung &
Prompt: Die Reihenfolge der members ist eigentlich egal: Ich möchte gerne eine Funktion schreiben, die für eine membersliste folgendes macht: Schablonartig übereinander legen und schauen, ob es Übereinstimmungen gibt. Erst basierend auf der Reihenfolge, aber unberücksichtigt eventuell zusätzlicher Wörter zwischen den Wörtern. Das sollte absichern, dass alle "Declaration xy by abc" gefunden werden mit einem Matcher "Declaration % by %" - Unübereinstimmungen sollen mit % gekennzeichnet werden. Für die beiden Payment-Actions kommt dann vermutlich "%Payment%" raus. Egal, kein Overfitting:
Ich glaube, dafür reicht entweder ein Textmatcher, aaaber es fühlt sich für mich auch noch Process Mining vom Vorgehen an, wo Ähnlichkeiten zweier Logs untersucht werden - vielleicht ist deswegen auch pm4py eine gute Möglichkeit. &
Claude (Anthropic), Opus 5.5 07.10.26 // 6.8 Credits (GitHub Copilot, Max/1M) \\
\hline
Syntaxhilfe und Styling von Diagrammen und technischer Projektstruktur &
Als Werkzeug zur Lösung von Syntaxproblemen, Inkonsistenzen in meinen Befehlen, aufgetretenden Refactoring-Problemen oder das Anpassen der Ausgabe-Styles von BPMN-Diagrammen oder Tabellen. Außerdem als Hilfe bei Problemen in der Programmumgebung. &
Aufbau, Reihenfolge und inhaltliche Festlegungen stammen von mir, die KI hat sie in Code übersetzt. Vorschläge wurden ausnahmslos als lokaler Git-Diff gesichtet und erst nach eigener Prüfung übernommen. &
Claude (Anthropic), Opus 4.8 (GitHub Copilot, Max/1M) \\
\hline
Umbenennung und Refaktorisierung der Jupyter-Notebooks &
Erhöhung der Übersicht und Vereinheitlichung der Projektstruktur durch Auslagerung in extra Funktionspackage &
Rein organisatorische Änderung, das Ergebnis habe ich gegen die Notebook-Reihenfolge geprüft. Ohne Einfluss auf den Inhalt der Arbeit. &
Claude (Anthropic), Opus 4.8 (GitHub Copilot, Max/1M) \\
\hline
\end{longtable}
\end{landscape}
